\section{Conclusion}
This paper investigates how the design of an audit standard can undermine its effectiveness, complicate efforts to assess audit compliance with the standard, and undermine accountability through a case study of ASB 018, a standard for auditing probabilistic genotyping software. Through qualitative analysis of ASB 018 and five publicly available audit reports, we identify numerous gaps between the standard's desired outcomes and the auditing practices it enables. We identify a number of design features that enable these gaps, such as the standard's failure to define key concepts, requirements that treat audit components as separate considerations that can be assessed independently of one another, and vague language. Based on our findings, we recommend that designers of audit standards clearly articulate their goals, carefully examine how designing for compatibility with users' practices can undermine effectiveness, ensure sufficient specificity in the requirements language to ensure effectiveness and accountability, and co-design standards with stakeholders beyond the standard's intended users. 